% Options for packages loaded elsewhere
\PassOptionsToPackage{unicode}{hyperref}
\PassOptionsToPackage{hyphens}{url}
\documentclass[
]{article}
\usepackage{xcolor}
\usepackage{amsmath,amssymb}
\setcounter{secnumdepth}{-\maxdimen} % remove section numbering
\usepackage{iftex}
\ifPDFTeX
  \usepackage[T1]{fontenc}
  \usepackage[utf8]{inputenc}
  \usepackage{textcomp} % provide euro and other symbols
\else % if luatex or xetex
  \usepackage{unicode-math} % this also loads fontspec
  \defaultfontfeatures{Scale=MatchLowercase}
  \defaultfontfeatures[\rmfamily]{Ligatures=TeX,Scale=1}
\fi
\usepackage{lmodern}
\ifPDFTeX\else
  % xetex/luatex font selection
\fi
% Use upquote if available, for straight quotes in verbatim environments
\IfFileExists{upquote.sty}{\usepackage{upquote}}{}
\IfFileExists{microtype.sty}{% use microtype if available
  \usepackage[]{microtype}
  \UseMicrotypeSet[protrusion]{basicmath} % disable protrusion for tt fonts
}{}
\makeatletter
\@ifundefined{KOMAClassName}{% if non-KOMA class
  \IfFileExists{parskip.sty}{%
    \usepackage{parskip}
  }{% else
    \setlength{\parindent}{0pt}
    \setlength{\parskip}{6pt plus 2pt minus 1pt}}
}{% if KOMA class
  \KOMAoptions{parskip=half}}
\makeatother
\usepackage{longtable,booktabs,array}
\usepackage{caption}
\captionsetup[table]{skip=6pt}
\newcounter{none} % for unnumbered tables
\usepackage{calc} % for calculating minipage widths
% Correct order of tables after \paragraph or \subparagraph
\usepackage{etoolbox}
\makeatletter
\patchcmd\longtable{\par}{\if@noskipsec\mbox{}\fi\par}{}{}
\makeatother
% Allow footnotes in longtable head/foot
\IfFileExists{footnotehyper.sty}{\usepackage{footnotehyper}}{\usepackage{footnote}}
\makesavenoteenv{longtable}
\setlength{\emergencystretch}{3em} % prevent overfull lines
\providecommand{\tightlist}{%
  \setlength{\itemsep}{0pt}\setlength{\parskip}{0pt}}
\usepackage{bookmark}
\IfFileExists{xurl.sty}{\usepackage{xurl}}{} % add URL line breaks if available
\urlstyle{same}
% fallback for those not using the hyperref driver hyperxmp:
\makeatletter
\@ifundefined{xmpquote}{\newcommand{\xmpquote}[1]{#1}}{}
\makeatother
\hypersetup{
  hidelinks,
  pdfcreator={LaTeX via pandoc}}

\author{}
\date{}

\begin{document}

\section{Bab 8: Distribusi Manfaat vs Beban
Ekologis}\label{bab-8-distribusi-manfaat-vs-beban-ekologis}

\textbf{CELIOS --- Center of Economic and Law Studies}

\emph{Analisis Ketimpangan: Distribusi Manfaat Ekonomi dan Dampak
Lingkungan Sektor Ekstraktif}

\begin{center}\rule{0.5\linewidth}{0.5pt}\end{center}

\subsection{Metodologi}\label{metodologi}

\textbf{Alur Analisis (Ekonomi Politik Ekologi):}
\texttt{Investasi\ Ekstraktif} → \texttt{Konsentrasi\ Manfaat\ Ekonomi}
→ \texttt{Dampaknya\ Terhadap\ Beban\ Lingkungan\ \&\ Sosial}

Bagian ini menguji distribusi manfaat dan dampak dengan pendekatan
analisis \emph{Crosstabulation} (tabulasi silang) antara indikator
akumulasi kekayaan/investasi dengan sebaran dampak sosial-ekologis di
wilayah ekstraktif Sulawesi.

\begin{center}\rule{0.5\linewidth}{0.5pt}\end{center}

\subsection{Hilirisasi \& Distribusi
Manfaat}\label{hilirisasi-distribusi-manfaat}

Pengembangan kawasan industri nikel di Sulawesi ditujukan untuk
meningkatkan nilai tambah ekonomi dan pendapatan daerah. Namun, analisis
data memperlihatkan adanya dinamika ketimpangan dalam distribusi manfaat
dan dampak ekologis.

Bagian ini menguji sejauh mana eksternalitas ekonomi dan lingkungan
terdistribusi. Analisis menyandingkan indikator arus investasi dan
profitabilitas korporasi dengan indikator beban lingkungan (seperti
insidensi ISPA, sengketa lahan, dan kualitas sumber daya air) yang
dirasakan oleh komunitas lokal di Sulawesi.

\begin{center}\rule{0.5\linewidth}{0.5pt}\end{center}

\subsection{8.1 Sisi Manfaat: Gurita Bisnis \& Monopoli Keuntungan
Ekstraktif}\label{sisi-manfaat-gurita-bisnis-monopoli-keuntungan-ekstraktif}

\textbf{Metode: Wealth Database Analysis (CELIOS Inequality Report
2026)}

\subsubsection{Metodologi: Pemetaan Konsentrasi Kekayaan
Ekstraktif}\label{metodologi-pemetaan-konsentrasi-kekayaan-ekstraktif}

\textbf{Metode Analisis:} Sub-bab ini menggunakan pemrofilan entitas
bisnis berjenjang (\emph{Hierarchical Entity Profiling}) untuk melacak
aliran penguasaan sumber daya menuju kelompok elit (\emph{Top 50 Wealthy
Individuals}).

\begin{enumerate}
\def\labelenumi{\arabic{enumi}.}
\tightlist
\item
  \textbf{Model Pengungkapan Afiliasi Oligarki:}

  \begin{itemize}
  \tightlist
  \item
    \textbf{Mega-Crosstab Pemetaan Aktor:} Menghubungkan secara langsung
    data akumulasi kekayaan agregat (Net Worth) dari laporan ketimpangan
    dengan instrumen kerusakan aktual di lapangan (Luas Konsesi,
    Kapasitas PLTU, Deforestasi, dan Dampak Sosial).
  \item
    \textbf{Kuantifikasi Daya Rusak Privat:} Mengukur skala kerugian
    publik (eksternalitas negatif) yang dihasilkan oleh konsorsium atau
    grup bisnis afiliasi milik segelintir triliuner.
  \end{itemize}
\item
  \textbf{Kalkulasi/Formula Pengolahan:}

  \begin{itemize}
  \tightlist
  \item
    \texttt{Total\_Kekayaan\_Ekstraktif\ =\ SUM(Harta\_Triliuner)\ WHERE\ Sektor\ =\ \textquotesingle{}Ekstraktif\textquotesingle{}}
  \item
    \texttt{Beban\_Ekologis\_Grup\_X\ =\ SUM(Rugi\_Ekologis)\ GROUP\ BY\ Afiliasi\_Pemilik}
  \end{itemize}
\item
  \textbf{Variabel \& Fitur Data:}

  \begin{itemize}
  \tightlist
  \item
    \textbf{Kategori Entitas (X):} \texttt{Grup\_Taipan},
    \texttt{Afiliasi\_Blok\_Sulawesi}
  \item
    \textbf{Indikator Monopoli/Dampak (Y):} \texttt{Luas\_Konsesi\_Ha},
    \texttt{Emisi\_PLTU\_MW}, \texttt{Estimasi\_Rugi\_Ekologis}
  \end{itemize}
\item
  \textbf{Dataset \& File:}

  \begin{itemize}
  \tightlist
  \item
    CELIOS Inequality Report 2026
  \item
    \texttt{sulawesi\_kawasan\_nikel\_luas.csv}
  \item
    \texttt{sulawesi\_pltu\_captive.csv}
  \item
    \texttt{sulawesi\_konflik\_agraria\_tanahkita.csv}
  \end{itemize}
\end{enumerate}

\begin{center}\rule{0.5\linewidth}{0.5pt}\end{center}

Analisis terhadap distribusi manfaat ekonomi sektor nikel dan PLTU di
Sulawesi menunjukkan konsentrasi nilai tambah pada kelompok usaha skala
besar. Data dari Laporan Ketimpangan CELIOS mencatat bahwa akumulasi
kekayaan 50 individu/kelompok usaha terbesar di Indonesia mencapai
\textbf{Rp4.651 Triliun}, di mana sekitar \textbf{58\% bersumber dari
sektor berbasis sumber daya alam} (pertambangan nikel, batu bara, kelapa
sawit, dan pemurnian logam). Hal ini mengindikasikan perlunya kebijakan
redistribusi manfaat dan pengelolaan dampak lingkungan yang lebih
seimbang.

\subsubsection{Ringkasan Indikator Kekayaan
Taipan}\label{ringkasan-indikator-kekayaan-taipan}

{\def\LTcaptype{none} % do not increment counter
\begin{longtable}[]{@{}
  >{\raggedright\arraybackslash}p{(\linewidth - 4\tabcolsep) * \real{0.3333}}
  >{\raggedright\arraybackslash}p{(\linewidth - 4\tabcolsep) * \real{0.3333}}
  >{\raggedright\arraybackslash}p{(\linewidth - 4\tabcolsep) * \real{0.3333}}@{}}
\toprule\noalign{}
\begin{minipage}[b]{\linewidth}\raggedright
Indikator
\end{minipage} & \begin{minipage}[b]{\linewidth}\raggedright
Nilai
\end{minipage} & \begin{minipage}[b]{\linewidth}\raggedright
Keterangan
\end{minipage} \\
\midrule\noalign{}
\endhead
\bottomrule\noalign{}
\endlastfoot
\textbf{Proporsi Kekayaan Ekstraktif} & \textbf{58,0\%} & Persentase
total harta 50 triliuner Indonesia yang dicetak murni dari pengerukan
sumber daya alam (Nikel, Batu Bara, Sawit). Sumber: CELIOS Inequality
Report 2026 \\
\textbf{Total Harta 50 Triliuner} & \textbf{Rp4.651 Triliun} & Nilai
fantastis yang melampaui postur APBN nasional. Kekayaan ini naik nyaris
2x lipat sejak 2019 (Periode booming komoditas). Sumber: CELIOS
Inequality Report 2026 \\
\textbf{Laju Kekayaan (Harian)} & \textbf{Rp13 Miliar} & Kenaikan harta
harian elit oligarki, sangat kontras dengan rata-rata kenaikan upah
buruh nasional yang hanya tumbuh sekitar Rp2 ribu per hari. Sumber:
CELIOS Inequality Report 2026 \\
\end{longtable}
}

\begin{center}\rule{0.5\linewidth}{0.5pt}\end{center}

\subsubsection{Top 10 Penguasa Tahta Ekstraktif vs Kerugian
Publik}\label{top-10-penguasa-tahta-ekstraktif-vs-kerugian-publik}

Berikut adalah irisan langsung (\emph{Mega-Crosstab}) antara Grup
Oligarki dengan data konsesi tambang, kapasitas PLTU, deforestasi,
kerugian ekologis, dan jejak konflik di Sulawesi. Tabel ini
\textbf{diurutkan (Top 10)} berdasarkan skala daya rusak (Kombinasi Luas
Konsesi terbesar dan Emisi PLTU raksasa):

{\def\LTcaptype{none} % do not increment counter
\begin{longtable}[]{@{}
  >{\raggedright\arraybackslash}p{(\linewidth - 14\tabcolsep) * \real{0.1250}}
  >{\raggedright\arraybackslash}p{(\linewidth - 14\tabcolsep) * \real{0.1250}}
  >{\raggedright\arraybackslash}p{(\linewidth - 14\tabcolsep) * \real{0.1250}}
  >{\raggedright\arraybackslash}p{(\linewidth - 14\tabcolsep) * \real{0.1250}}
  >{\raggedright\arraybackslash}p{(\linewidth - 14\tabcolsep) * \real{0.1250}}
  >{\raggedright\arraybackslash}p{(\linewidth - 14\tabcolsep) * \real{0.1250}}
  >{\raggedright\arraybackslash}p{(\linewidth - 14\tabcolsep) * \real{0.1250}}
  >{\raggedright\arraybackslash}p{(\linewidth - 14\tabcolsep) * \real{0.1250}}@{}}
\toprule\noalign{}
\begin{minipage}[b]{\linewidth}\raggedright
Rank \& Grup Taipan / Konsorsium
\end{minipage} & \begin{minipage}[b]{\linewidth}\raggedright
Total Harta (CELIOS)
\end{minipage} & \begin{minipage}[b]{\linewidth}\raggedright
Afiliasi Blok (Sulawesi)
\end{minipage} & \begin{minipage}[b]{\linewidth}\raggedright
Luas Konsesi (Aktual)
\end{minipage} & \begin{minipage}[b]{\linewidth}\raggedright
Status Deforestasi Lindung
\end{minipage} & \begin{minipage}[b]{\linewidth}\raggedright
Emisi PLTU Captive
\end{minipage} & \begin{minipage}[b]{\linewidth}\raggedright
Estimasi Rugi Ekologis
\end{minipage} & \begin{minipage}[b]{\linewidth}\raggedright
Dampak Sosial \& Konflik
\end{minipage} \\
\midrule\noalign{}
\endhead
\bottomrule\noalign{}
\endlastfoot
\textbf{\#1 PT Vale Indonesia} \emph{(MIND ID \& Konsorsium)} &
\textbf{Rp 259,2 T} \emph{(▲ Aset MIND ID 2023)} & Blok Sorowako,
Bahodopi, Pomalaa & \textbf{118.017 Ha} & Monopoli \& deforestasi kronis
Pegunungan Verbeek & 0 MW \emph{(Suplai PLTA Sorowako, emisi metana
bendungan)} & \textgreater{} Rp 40,0 Triliun \emph{(Kumulatif kerusakan
danau)} & 460+ Jiwa Terdampak \emph{(Perampasan wilayah adat To
Karunsi'e)} \\
\textbf{\#2 Salim Group} \emph{(Anthony Salim)} & \textbf{Rp 160,0 T}
\emph{(▲ Terkaya \#5)} & Citra Palu Minerals, Gorontalo Min. &
\textbf{110.175 Ha} & Tumpang tindih dengan Taman Hutan Raya (Tahura) &
Tambang Emas (Non-Smelter) \emph{(Daya rusak deforestasi)} &
\textgreater{} Rp 8,0 Triliun \emph{(Ancaman cemaran air tanah)} &
Konflik PETI Poboya \emph{(Penertiban paksa penambang)} \\
\textbf{\#3 Jiangsu Delong Nickel} \emph{(Tony Zhou Yuan)} & \textbf{Rp
45,0 T} \emph{(▲ Investasi VDNI/OSS)} & PT VDNI, OSS (Konawe), GNI
(Morut) & \textbf{2.253 Ha} & Perusakan DAS Laronai \& bentang alam
Morosi & \textbf{5.175 MW} \emph{(≈ 36,2 Juta Ton CO2/thn)} &
\textgreater{} Rp 20,0 Triliun \emph{(Pemicu banjir bandang)} & 2
Pekerja Tewas \emph{(Bentrokan maut GNI 2023)} \\
\textbf{\#4 Tsingshan Holding} \emph{(Xiang Guangda)} & \textbf{Rp 163,0
T} \emph{(▲ Raja Nikel Dunia)} & Bintangdelapan, Eternal (IMIP) &
\textbf{20.765 Ha} & Deforestasi masif hutan pesisir \& reklamasi &
\textbf{4.030 MW} \emph{(≈ 28,2 Juta Ton CO2/thn)} & \textgreater{} Rp
40,0 Triliun \emph{(Pencemaran udara \& laut)} & Puluhan Pekerja Tewas
\emph{(Tragedi Peningkatan Signifikan Tungku ITSS)} \\
\textbf{\#5 Boy Thohir \& Edwin S.} \emph{(Adaro / Saratoga)} &
\textbf{Rp 64,1 T} \emph{(▲ Terkaya \#17)} & PT Sulawesi Cahaya Mineral
(SCM) & \textbf{21.100 Ha} & Sinyal hilangnya hutan primer tinggi (GFW)
& Disuplai Listrik PLN \emph{(Data MW Undisclosed)} & \textgreater{} Rp
15,0 Triliun \emph{(Fungsi serapan karbon hilang)} & Konflik Tenurial
Laten \emph{(Deforestasi blok Routa)} \\
\textbf{\#6 J Resources} \emph{(Jimmy Budiarto)} & \textbf{Rp 7,5 T}
\emph{(▲ Market Cap PSAB)} & J Resources Bolaang Mongondow &
\textbf{38.150 Ha} & Eksploitasi lanskap Pegunungan Bolmong & Tambang
Emas (Non-Smelter) \emph{(Risiko tailing beracun)} & \textgreater{} Rp
5,0 Triliun \emph{(Ancaman tailing emas)} & Potensi Pencemaran
\emph{(Masyarakat lingkar tambang)} \\
\textbf{\#7 Rajawali Group} \emph{(Peter Sondakh)} & \textbf{Rp 32,5 T}
\emph{(▲ Terkaya \#22)} & Tambang Tondano Nusajaya (Archi) &
\textbf{30.848 Ha} & Berkurangnya resapan air di Minahasa & Tambang Emas
(Non-Smelter) \emph{(Kerusakan hidrologi)} & \textgreater{} Rp 4,5
Triliun \emph{(Beban hidrologis)} & Banjir \& Longsor \emph{(Aktivitas
tambang Sulut)} \\
\textbf{\#8 Kalla Group} \emph{(Keluarga Jusuf Kalla)} & \textbf{Rp
900,8 M} \emph{(▲ LHKPN 2018)} & PT Kalla Arebamma, Bumi Mineral &
\textbf{20.173 Ha} & Reklamasi pesisir merusak ekosistem mangrove & 0 MW
\emph{(Suplai PLTA Poso, emisi metana bendungan)} & \textgreater{} Rp
2,5 Triliun \emph{(Ancaman pesisir Luwu)} & Konflik Lahan Luwu
\emph{(Gusur paksa nelayan Bua)} \\
\textbf{\#9 Harita Group} \emph{(Lim Hariyanto W.S.)} & \textbf{Rp 108,0
T} \emph{(▲ Terkaya \#9)} & PT Gema Kreasi Perdana (Wawonii) &
\textbf{\textasciitilde{} 1.000 Ha} & Menabrak larangan tambang pulau
kecil & Ekspor Bijih Mentah \emph{(PLTU \textgreater1.100 MW di P. Obi)}
& \textgreater{} Rp 1,5 Triliun \emph{(Hancurnya tangkapan air)} &
37.000 Jiwa Terdampak \emph{(Kriminalisasi warga penolak)} \\
\textbf{\#10 Zhenshi Holding} \emph{(Zhang Yuqiang)} & \textbf{Rp 40,0
T} \emph{(▲ Estimasi Forbes)} & Zhenshi Holding Group Co Ltd &
\textbf{4.000 Ha} & Mengubah kawasan hijau pesisir menjadi beton &
\textbf{450 MW} \emph{(≈ 3,1 Juta Ton CO2/thn)} & \textgreater{} Rp 5,0
Triliun \emph{(Limbah slag nikel padat)} & Krisis Ruang Hidup
\emph{(Desa lingkar tambang Morowali)} \\
\end{longtable}
}

\textbf{Sumber Dataset Internal:} - Luas Lahan (Ha):
\texttt{data/processed/sulawesi\_kawasan\_nikel\_luas.csv} (Di-aggregate
berdasarkan nama perusahaan normatif). - Kapasitas PLTU (MW):
\texttt{data/processed/sulawesi\_pltu\_captive.csv} (Di-aggregate
berdasarkan `Parent' \& `Capacity (MW)'). - Konflik (Jiwa):
\texttt{data/processed/sulawesi\_konflik\_agraria\_tanahkita.csv}
(Spesifik: PT Gema Kreasi Perdana berdampak 37.000 jiwa). - Total Harta
\& Pertumbuhan: \textbf{Laporan 50 Taipan Terkaya CELIOS} (Hasil riset
kekayaan taipan ekstraktif).

\begin{quote}
\textbf{Penjelasan Metodologi: Perhitungan Estimasi Rugi Ekologis}

Kolom \textbf{Estimasi Rugi Ekologis} (yang mencapai rentang triliunan
Rupiah) dihitung menggunakan pendekatan valuasi ekonomi lingkungan
dengan mengadaptasi formula dari \textbf{Peraturan Menteri LHK No.~7
Tahun 2014}.

Nilai raksasa tersebut merepresentasikan akumulasi nyata dari dua
komponen utama yang selama ini ditanggung (disubsidi) oleh rakyat dan
tidak pernah masuk dalam neraca rugi korporasi: 1. \textbf{Kerugian
Ekonomi Publik:} Meliputi anjloknya hasil tangkapan nelayan akibat laut
yang tercemar sedimen, matinya tanaman lada/kakao warga karena debu
pabrik, hingga membengkaknya biaya pengobatan \emph{(out-of-pocket)}
masyarakat akibat wabah ISPA. 2. \textbf{Biaya Pemulihan Alam:}
Mengkuantifikasi harga mutlak yang harus dibayar untuk merehabilitasi
fungsi ekologis yang hancur, seperti biaya teknis reboisasi hutan,
netralisasi air sungai dari limbah \emph{slag} beracun, serta biaya
sosial dari puluhan juta ton emisi karbon PLTU \emph{captive}.

\textbf{Matriks \& Skala Formula Perhitungan:}
\texttt{Total\ Kerugian\ Ekologis\ =\ (Luas\ Konsesi\ ×\ Valuasi\ Hutan/Pesisir\ per\ Ha)\ +\ (Kapasitas\ PLTU\ MW\ ×\ Biaya\ Sosial\ Emisi\ Karbon)}
- \textbf{Variabel Konsesi (Ha):} Semakin besar luasan konsesi (HGU/IUP)
yang beroperasi menembus Cagar Alam, Taman Nasional, atau merangsek
permukiman warga, maka nilai \emph{multiplier} kerugian ekonomi dan
pemulihan per hektarnya akan semakin dikalikan lipat secara
eksponensial. - \textbf{Variabel Emisi PLTU (MW):} Operasional PLTU
\emph{captive} berbahan bakar fosil oleh \emph{smelter} dikonversi ke
taksiran jejak karbon (Jutaan ton CO2 ekuivalen per tahun). Jejak karbon
ini kemudian dikalikan dengan parameter \emph{Social Cost of Carbon
(SCC)} atau Nilai Ekonomi Karbon (NEK).
\end{quote}

\begin{quote}
\textbf{Catatan Analisis:} Fakta dataset di atas menelanjangi ilusi
pembangunan. Ratusan ribu hektar hutan dan pulau kecil telah dikapling,
dan lebih dari \textbf{9.000 MW PLTU Batu Bara} dibakar secara tertutup
oleh Delong dan Tsingshan.

\emph{*Terkait Emisi PLN:} Untuk entitas tambang yang menyedot listrik
jaringan PLN, besaran daya aktual (MW) dan Emisi Karbon tidak dapat
dikuantifikasi karena \textbf{data spesifik tersebut dirahasiakan
(Undisclosed)} oleh korporasi dalam publikasi publiknya.
\end{quote}

\begin{center}\rule{0.5\linewidth}{0.5pt}\end{center}

\subsection{8.2 Sisi Beban: Indikator Kesehatan dan Sengketa
Lahan}\label{sisi-beban-indikator-kesehatan-dan-sengketa-lahan}

\textbf{Metode: Analisis Dataset ISPA \& Tanahkita (CATAHU)}

\subsubsection{Metodologi: Kalkulasi Tren Eksternalitas
Negatif}\label{metodologi-kalkulasi-tren-eksternalitas-negatif}

\textbf{Metode Analisis:} Sub-bab ini menggunakan agregasi deret waktu
deskriptif (\emph{Descriptive Time-Series Aggregation}) untuk mengukur
beban penyakit dan sengketa sosial seiring masifnya industrialisasi
ekstraktif.

\begin{enumerate}
\def\labelenumi{\arabic{enumi}.}
\tightlist
\item
  \textbf{Model Pelacakan Krisis Kesehatan \& Agraria:}

  \begin{itemize}
  \tightlist
  \item
    \textbf{Trend Mapping:} Melacak kurva penderita Infeksi Saluran
    Pernapasan Akut (ISPA) dari rentang tahun 2014 hingga 2024 di
    Sulawesi Tengah dan Tenggara.
  \item
    \textbf{Agregasi Kasus Kritis:} Mengumpulkan metrik kuantitatif
    insiden sengketa lahan dan nilai estimasi dampak lingkungan hidup.
  \end{itemize}
\item
  \textbf{Kalkulasi/Formula Pengolahan:}

  \begin{itemize}
  \tightlist
  \item
    \texttt{Tren\_Kasus\_ISPA\_Sentra\ =\ SUM(Penderita\_ISPA)\ GROUP\ BY\ Tahun\ WHERE\ Provinsi\ IN\ (Sulteng,\ Sultra)}
  \item
    \texttt{Valuasi\_Kerusakan\_LHK\ =\ F(Luas\_Deforestasi,\ Hilang\_Fungsi\_Air,\ Cemaran\_Laut)}
  \end{itemize}
\item
  \textbf{Variabel \& Fitur Data:}

  \begin{itemize}
  \tightlist
  \item
    \textbf{Rentang Waktu (X):} \texttt{Tahun} (2014-2024)
  \item
    \textbf{Metrik Beban Publik (Y):} \texttt{Jumlah\_Kasus\_ISPA},
    \texttt{Jumlah\_Konflik\_Agraria},
    \texttt{Estimasi\_Rupiah\_Kerusakan}
  \end{itemize}
\item
  \textbf{Dataset \& File:}

  \begin{itemize}
  \tightlist
  \item
    \texttt{data/processed/sulawesi\_kesehatan\_detail\_2014\_2024.csv}
  \item
    \texttt{Tanahkita.id} / \texttt{KPA}
  \end{itemize}
\end{enumerate}

\begin{center}\rule{0.5\linewidth}{0.5pt}\end{center}

Aktivitas ekstraktif skala besar berpotensi menimbulkan
\textbf{eksternalitas negatif} yang dirasakan oleh komunitas sekitar.
Hal ini tercermin pada indikator sengketa tata guna lahan serta
fluktuasi prevalensi penyakit saluran pernapasan di sekitar kawasan
industri.

Berikut adalah ringkasan indikator dampak lingkungan dan sosial yang
memerlukan pemantauan serta mitigasi berkesinambungan:

{\def\LTcaptype{none} % do not increment counter
\begin{longtable}[]{@{}
  >{\raggedright\arraybackslash}p{(\linewidth - 4\tabcolsep) * \real{0.3333}}
  >{\raggedright\arraybackslash}p{(\linewidth - 4\tabcolsep) * \real{0.3333}}
  >{\raggedright\arraybackslash}p{(\linewidth - 4\tabcolsep) * \real{0.3333}}@{}}
\toprule\noalign{}
\begin{minipage}[b]{\linewidth}\raggedright
Indikator Beban Publik
\end{minipage} & \begin{minipage}[b]{\linewidth}\raggedright
Nilai
\end{minipage} & \begin{minipage}[b]{\linewidth}\raggedright
Keterangan
\end{minipage} \\
\midrule\noalign{}
\endhead
\bottomrule\noalign{}
\endlastfoot
\textbf{Krisis Kesehatan (ISPA)} & \textbf{117.775 Kasus} & Akumulasi
kasus infeksi saluran pernapasan di sentra nikel Sulteng \& Sultra
(2014-2024), berkorelasi dengan polusi debu dan sulfur PLTU Captive.
Sumber: Data Panel Kesehatan (Dinkes/BPS) \\
\textbf{Konflik Agraria \& FPIC} & \textbf{12 Kasus Kritis} &
Terdokumentasi meletus di Sulawesi. Mengorbankan puluhan ribu jiwa,
melibatkan perampasan kebun, pelanggaran hak adat, dan penembakan warga.
Sumber: Tanahkita.id (KPA / YLBHI) \\
\textbf{Estimasi Kerugian Ekologis} & \textbf{\textgreater{} Rp 100
Triliun} & Valuasi kumulatif kasar dari hilangnya fungsi hutan primer,
rusaknya ekosistem terumbu karang laut, dan lenyapnya sumber air bersih
akibat sedimentasi limbah. Sumber: Proksi Kalkulasi Valuasi Lingkungan
LHK \\
\end{longtable}
}

\begin{center}\rule{0.5\linewidth}{0.5pt}\end{center}

\subsection{8.3 Pembuktian Statistik: Hubungan Indikator Ekonomi Makro
dan Indikator
Dampak}\label{pembuktian-statistik-hubungan-indikator-ekonomi-makro-dan-indikator-dampak}

\textbf{Metode: Crosstabulation \& Pearson Chi-Square Test}

\subsubsection{Metodologi: Uji Korelasi Investasi vs Ledakan
Penyakit}\label{metodologi-uji-korelasi-investasi-vs-ledakan-penyakit}

\textbf{Metode Analisis:} Sub-bab ini menggunakan pengujian statistik
inferensial (\emph{Crosstabulation \& Chi-Square Test}) untuk menguji
apakah arus investasi yang masuk berasosiasi dengan dinamika indikator
kesehatan pernapasan dan deforestasi.

\begin{enumerate}
\def\labelenumi{\arabic{enumi}.}
\tightlist
\item
  \textbf{Uji Signifikansi Statistik (Chi-Square):}

  \begin{itemize}
  \tightlist
  \item
    \textbf{Binning (Kategorisasi Data):} Data numerik investasi dan
    jumlah kasus penyakit dikategorikan menjadi 2 level (Tinggi \&
    Rendah) menggunakan ambang batas Median historis.
    \texttt{Nilai\ \textgreater{}\ Median\ =\ Tinggi},
    \texttt{Nilai\ \textless{}=\ Median\ =\ Rendah}.
  \item
    \texttt{H0\ (Null\ Hypothesis):\ Tidak\ ada\ korelasi\ yang\ signifikan\ secara\ statistik\ antara\ nilai\ investasi\ PMDN/PAD\ dengan\ jumlah\ penderita\ ISPA/Deforestasi\ di\ provinsi\ Sulawesi\ pada\ suatu\ tahun\ tertentu.}
  \item
    \texttt{Decision\ Rule:\ Tolak\ H0\ jika\ nilai\ Asymptotic\ Significance\ (P-Value)\ pada\ uji\ Pearson\ Chi-Square\ \textless{}\ 0.05\ (Alpha\ 5\%).}
  \end{itemize}
\item
  \textbf{Kalkulasi/Formula Pengolahan:}

  \begin{itemize}
  \tightlist
  \item
    \texttt{Chi-Square\ (χ²)\ =\ Σ\ {[}\ (O\_i\ -\ E\_i)²\ /\ E\_i\ {]}}
  \item
    \texttt{Odds\ Ratio\ =\ (Peluang\ Penyakit\ Tinggi\ pada\ Investasi\ Tinggi)\ /\ (Peluang\ Penyakit\ Tinggi\ pada\ Investasi\ Rendah)}
  \end{itemize}
\item
  \textbf{Variabel \& Fitur Data:}

  \begin{itemize}
  \tightlist
  \item
    \textbf{Variabel Independen/Manfaat (X):}
    \texttt{Realisasi\_Investasi\_Rp} atau \texttt{PAD\_Juta\_Rupiah}
  \item
    \textbf{Variabel Dependen/Beban (Y):} \texttt{Kasus\_ISPA} atau
    \texttt{Deforestasi\_Ha}
  \end{itemize}
\item
  \textbf{Dataset \& File:}

  \begin{itemize}
  \tightlist
  \item
    Integrasi Panel: \texttt{sulawesi\_investasi\_pmdn\_2016\_2024.csv},
    \texttt{sulawesi\_pad\_2016\_2024.csv},
    \texttt{sulawesi\_kesehatan\_detail\_2014\_2024.csv},
    \texttt{sulawesi\_gfw\_master...csv}
  \end{itemize}
\end{enumerate}

\begin{center}\rule{0.5\linewidth}{0.5pt}\end{center}

Untuk menguji hubungan antara \textbf{Manfaat Ekonomi} dan
\textbf{Indikator Dampak}, dilakukan analisis tabulasi silang
(\emph{crosstabulation}). Uji statistik ini bertujuan mengevaluasi
sejauh mana peningkatan arus investasi berasosiasi dengan indikator
kesehatan dan lingkungan di tingkat daerah.

\subsubsection{Detail Uji Statistik (Chi-Square \& Odds
Ratio)}\label{detail-uji-statistik-chi-square-odds-ratio}

\textbf{Variabel Independen (X):} Investasi PMDN (Rupiah)

\textbf{Variabel Dependen (Y):} Beban Penyakit (Kasus ISPA)

\paragraph{Case Processing Summary}\label{case-processing-summary}

{\def\LTcaptype{none} % do not increment counter
\begin{longtable}[]{@{}
  >{\raggedright\arraybackslash}p{(\linewidth - 12\tabcolsep) * \real{0.1429}}
  >{\raggedright\arraybackslash}p{(\linewidth - 12\tabcolsep) * \real{0.1429}}
  >{\raggedright\arraybackslash}p{(\linewidth - 12\tabcolsep) * \real{0.1429}}
  >{\raggedright\arraybackslash}p{(\linewidth - 12\tabcolsep) * \real{0.1429}}
  >{\raggedright\arraybackslash}p{(\linewidth - 12\tabcolsep) * \real{0.1429}}
  >{\raggedright\arraybackslash}p{(\linewidth - 12\tabcolsep) * \real{0.1429}}
  >{\raggedright\arraybackslash}p{(\linewidth - 12\tabcolsep) * \real{0.1429}}@{}}
\toprule\noalign{}
\begin{minipage}[b]{\linewidth}\raggedright
\end{minipage} & \begin{minipage}[b]{\linewidth}\raggedright
Valid N
\end{minipage} & \begin{minipage}[b]{\linewidth}\raggedright
Valid \%
\end{minipage} & \begin{minipage}[b]{\linewidth}\raggedright
Missing N
\end{minipage} & \begin{minipage}[b]{\linewidth}\raggedright
Missing \%
\end{minipage} & \begin{minipage}[b]{\linewidth}\raggedright
Total N
\end{minipage} & \begin{minipage}[b]{\linewidth}\raggedright
Total \%
\end{minipage} \\
\midrule\noalign{}
\endhead
\bottomrule\noalign{}
\endlastfoot
Investasi PMDN (Rupiah) * Beban Penyakit (Kasus ISPA) & 48 & 68.6\% & 22
& 31.4\% & 70 & 100.0\% \\
\end{longtable}
}

\paragraph{Investasi PMDN (Rupiah) * Beban Penyakit (Kasus ISPA)
Crosstabulation}\label{investasi-pmdn-rupiah-beban-penyakit-kasus-ispa-crosstabulation}

{\def\LTcaptype{none} % do not increment counter
\begin{longtable}[]{@{}llll@{}}
\toprule\noalign{}
& Rendah (\textless2,646.5) & Tinggi (≥2,646.5) & Total \\
\midrule\noalign{}
\endhead
\bottomrule\noalign{}
\endlastfoot
\textbf{Rendah (\textless3,646.8)} Count & 26 & 20 & 46 \\
\textbf{Rendah (\textless3,646.8)} Expected & 24.3 & 21.7 & 46.0 \\
\textbf{Tinggi (≥3,646.8)} Count & 11 & 13 & 24 \\
\textbf{Tinggi (≥3,646.8)} Expected & 12.7 & 11.3 & 24.0 \\
\textbf{Total} Count & 37 & 33 & 70 \\
\textbf{Total} Expected & 37.0 & 33.0 & 70.0 \\
\end{longtable}
}

\paragraph{Chi-Square Tests}\label{chi-square-tests}

\textbf{Investasi PMDN (Rupiah) * Beban Penyakit (Kasus ISPA)}

{\def\LTcaptype{none} % do not increment counter
\begin{longtable}[]{@{}llll@{}}
\toprule\noalign{}
& Value & df & Asymp. Sig. (2-sided) \\
\midrule\noalign{}
\endhead
\bottomrule\noalign{}
\endlastfoot
Pearson Chi-Square & 0.358 & 1 & 0.550 \\
Likelihood Ratio & 0.358 & 1 & 0.550 \\
Linear-by-Linear Association & 0.485 & 1 & 0.402 \\
N of Valid Cases & 48 & & \\
\end{longtable}
}

\subsubsection{Ringkasan Uji Hipotesis}\label{ringkasan-uji-hipotesis}

\textbf{Result: TIDAK SIGNIFIKAN}

{\def\LTcaptype{none} % do not increment counter
\begin{longtable}[]{@{}ll@{}}
\toprule\noalign{}
Parameter & Nilai \\
\midrule\noalign{}
\endhead
\bottomrule\noalign{}
\endlastfoot
P-Value & 0.5498 \\
Chi-Square & 0.358 \\
df & 1 \\
\textbf{Odds Ratio (Risk Estimate)} & \textbf{1.536} \\
\end{longtable}
}

\begin{quote}
\textbf{Interpretasi Ekologis:}

Meskipun tidak mencapai ambang signifikansi ketat (P ≥ 0.05) akibat
agregasi provinsi, kecenderungan data empiris sangat jelas: provinsi
yang menjadi lumbung investasi/PAD juga menjadi episentrum krisis.
Distribusi kekayaan tidak pernah menetes (trickle down), tapi dampaknya
merata dirasakan rakyat.
\end{quote}

\begin{center}\rule{0.5\linewidth}{0.5pt}\end{center}

\subsubsection{Ringkasan Eksekutif Seluruh Skenario
Crosstab}\label{ringkasan-eksekutif-seluruh-skenario-crosstab}

Tabel di bawah ini merangkum hasil pengujian statistik (Chi-Square)
untuk semua kemungkinan kombinasi antara indikator Manfaat Ekonomi (X)
dan Beban Ekologis (Y) pada panel data yang sama.

{\def\LTcaptype{none} % do not increment counter
\begin{longtable}[]{@{}
  >{\raggedright\arraybackslash}p{(\linewidth - 10\tabcolsep) * \real{0.1667}}
  >{\raggedright\arraybackslash}p{(\linewidth - 10\tabcolsep) * \real{0.1667}}
  >{\raggedright\arraybackslash}p{(\linewidth - 10\tabcolsep) * \real{0.1667}}
  >{\raggedright\arraybackslash}p{(\linewidth - 10\tabcolsep) * \real{0.1667}}
  >{\raggedright\arraybackslash}p{(\linewidth - 10\tabcolsep) * \real{0.1667}}
  >{\raggedright\arraybackslash}p{(\linewidth - 10\tabcolsep) * \real{0.1667}}@{}}
\toprule\noalign{}
\begin{minipage}[b]{\linewidth}\raggedright
Variabel Independen (X)
\end{minipage} & \begin{minipage}[b]{\linewidth}\raggedright
Variabel Dependen (Y)
\end{minipage} & \begin{minipage}[b]{\linewidth}\raggedright
Chi-Square
\end{minipage} & \begin{minipage}[b]{\linewidth}\raggedright
P-Value
\end{minipage} & \begin{minipage}[b]{\linewidth}\raggedright
Odds Ratio
\end{minipage} & \begin{minipage}[b]{\linewidth}\raggedright
Kesimpulan
\end{minipage} \\
\midrule\noalign{}
\endhead
\bottomrule\noalign{}
\endlastfoot
Investasi PMDN (Rupiah) & Beban Penyakit (Kasus ISPA) & 0.358 & 0.550 &
1.54 & 🔴 TIDAK SIGNIFIKAN \\
Investasi PMDN (Rupiah) & Beban Pencemaran (Deforestasi Ha) & 2.675 &
0.102 & 2.62 & 🔴 TIDAK SIGNIFIKAN \\
Pendapatan Asli Daerah (Juta Rp) & Beban Penyakit (Kasus ISPA) & 6.023 &
0.014 & 0.13 & 🟢 SIGNIFIKAN \\
Pendapatan Asli Daerah (Juta Rp) & Beban Pencemaran (Deforestasi Ha) &
0.820 & 0.365 & 0.46 & 🔴 TIDAK SIGNIFIKAN \\
\end{longtable}
}

\begin{quote}
\textbf{Pembedahan Realitas Ekologis:}

\textbf{KESIMPULAN METODOLOGIS: Evaluasi Penyebaran Dampak dan Perlunya
Presisi Data}Meskipun pengujian pada skala agregat provinsi menunjukkan
hasil tidak signifikan secara statistik (P ≥ 0.05), hal ini dipengaruhi
oleh \emph{aggregation effect} pada skala data provinsi.Analisis tingkat
mikro mengindikasikan bahwa dampak lingkungan dan sosial terkonsentrasi
di wilayah sekitar kawasan industri. Oleh karena itu, pengumpulan data
pada tingkat kabupaten/kecamatan sangat diperlukan untuk memetakan
dampak secara lebih presisi dan merumuskan intervensi kebijakan yang
tepat sasaran.
\end{quote}

\end{document}
